% from main.tex
%%%%%%%%%%  scan idx 538  |  folio 525  %%%%%%%%%%
Il s'%
ensuit (V) que $x$ est la $m$-ème fonction symétrique élémentaire en les $\cl^\circ(L_r)-1$ ($r=1,\ldots,n$), donc est un polynôme homogène de degré $m$ à coefficients entiers dans les variables $(\cl^\circ(L_r)-1)$ ($r=1,\ldots,n$). Appliquant à nouveau l'hypothèse de récurrence et l'associativité de la loi (1.3.1) on est ainsi ramené à démontrer que, pour tout $\mathcal O_X$-module inversible $L$, on a
\begin{equation}\tag{1.3.6}
(\cl^\circ(L)-1)\cl_\bullet(\mathcal O_X)\in\Filt_{j-1}(X)
\end{equation}
Le schéma $X$ étant intègre, $L$ et $\mathcal O_X$ sont isomorphes à des sous-modules du $\mathcal O_X$-module constant $R_X$ des fonctions rationnelles sur $X$. Le $\mathcal O_X$-module somme $F=\mathcal O_X+L$ est cohérent, et on trouve
\begin{align*}
(\cl^\circ(L)-1)\cl_\bullet(\mathcal O_X)
&=\cl_\bullet(L)-\cl_\bullet(\mathcal O_X)\\
&=(\cl_\bullet(F)-\cl_\bullet(\mathcal O_X))-(\cl_\bullet(F)-\cl_\bullet(L))\\
&=\cl_\bullet(F/\mathcal O_X)-\cl_\bullet(F/L).
\end{align*}
Mais les fibres au point générique des faisceaux $F/\mathcal O_X$ et $F/L$ sont nulles, donc leurs supports dans $X$ sont rares ($X$ étant irréductible), donc de dimension $\leq j-1$. L'assertion (1.3.6) en découle immédiatement.

II. Reste à démontrer le cas général. On introduit le fibré en drapeaux
\[
g:D\longrightarrow X
\]
de $E$ sur $X$, où on peut trouver des $\mathcal O_D$-modules inversibles $L_1,\ldots,L_n$ tels que
\begin{equation}\tag{1.3.7}
g^*(\cl^\circ(E))=\cl^\circ(L_1)+\cdots+\cl^\circ(L_n)\in K^\circ(D).
\end{equation}
Soient $e$ le point générique de $X$, $k(e)$ le corps résiduel de $e$. Alors on peut trouver un point rationnel $u$ de la fibre générique $D_e$ de $D$ sur $k(e)$, car ces
